\documentclass[10pt,a4paper]{amsart}
\usepackage[margin=19mm]{geometry}
\usepackage{amsmath,amssymb,mathtools}
\usepackage[scr=rsfs,cal=euler]{mathalfa}
\usepackage{xfrac}
\usepackage[unicode,hidelinks,bookmarks=false]{hyperref}
\newcommand{\ie}{i.e.\ }
\newcommand{\cf}{cf.\ }
\newcommand{\resp}{respectively\ }
\newcommand{\Subsection}{\subsection}
\providecommand{\nobreakdash}{\nobreak\hbox{-}}
\setlength{\emergencystretch}{2em}
\multlinegap0pt
\usepackage[shortlabels]{enumitem}
\newtheorem{assumption}{Assumption}
\newtheorem{lem}{Lemma}
\newtheorem{thm}{Theorem}
\newcommand{\E}{\mathbb{E}}
\newcommand{\R}{\mathbb{R}}
\newcommand{\dx}[1][x]{%
	\ensuremath{
		{\mathrm{d}}{#1}
	}
}
\newcommand{\hr}{\le_{\mathrm{hr}}}
\providecommand{\lc}{\le_{\mathrm{lc}}}
\newcommand{\lr}{\le_{\mathrm{lr}}}
\newcommand{\st}{\le_{\mathrm{st}}}
\title{Kernel Characterisations of Stochastic Orders Within Parametric Density Families}
\author{Zakaria Derbazi}
\date{Source: arXiv:2605.18751v1 (CC BY 4.0)}
\begin{document}
\maketitle

\noindent\emph{Source and licence.} Kernel Characterisations of Stochastic Orders Within Parametric Density Families. Version arXiv:2605.18751v1 (CC BY 4.0), distributed by arXiv under the Creative Commons Attribution 4.0 International licence, http://creativecommons.org/licenses/by/4.0/, as recorded in the arXiv licence metadata for this exact version and shown on the abstract page https://arxiv.org/abs/2605.18751v1. Full licence notice: \texttt{licenses/arxiv-2605.18751-CC-BY-4.0.txt}.\par\smallskip

\noindent\textit{Editorial scope.} Counted unit: the article's thm thm:four-criteria (source0.tex lines 267--279). The article's complete original proof of it is delivered with it (source0.tex lines 280--319, one \texttt{proof} environment of 40 lines). No mathematics is written, repaired or re-derived: the statement and the proof are the article's own lines, in the article's own notation, and no retained line is substituted or re-flowed.  Retained uncounted as the source-local context the proof needs: lines 102--109; lines 170--172; lines 235--243; lines 258--260.  Nothing was omitted from any retained span, so the package carries no omission record.  No retained line contains a citation, so the wrapper carries no bibliography and no bibliography entry.  No retained line contains a figure environment, an image-inclusion command or drawing code, so no image is delivered and the package declares no asset dependency.  Editorial formatting only, in addition: an AMS \texttt{amsart} preamble that restates the article's own declarations and macros used by the retained lines, namely the environments \texttt{assumption}, \texttt{lem}, \texttt{thm} on separate counters, together with the standard packages those lines need.  The retained environments print in this document as Assumption 1, Lemma 1, Theorem 1 in order of appearance, whereas the article prints the same items with the numbering of its own full document.  The complete native source of the article as distributed in the arXiv e-print for this exact version is bundled unchanged as \texttt{sources/200873/source0.tex}.
\begin{assumption}[Differentiation under the integral]
\label{asm:diff-under-integral}
For every parametric density family \((f_\nu)_{\nu\in\mathcal I}\) considered below, with associated probability measures \((P_\nu)_{\nu\in\mathcal I}\) and \(\mathcal I\subseteq\R\) an interval, the density \(f_\nu(x)\) is positive on a common interval \(J\subseteq E\) and \(\nu \mapsto\log f_\nu(x)\) is \(C^1\) on \(\mathcal I\) for every \(x\in J\). For every compact interval \(B\subseteq\mathcal I\), there exists a majorant \(m\in L^1(J,\mu)\) such that
\[
\sup_{\nu\in B}\bigl|\partial_\nu f_\nu(x)\bigr|\le m(x)
\qquad \mu\text{-a.e.}
\]
\end{assumption}

\begin{equation}\label{def.log.companion}
s_\nu(x)=K_\nu(x)-\int_J K_\nu(y)\,\dx[]P_\nu(y),\qquad x\in J.
\end{equation}

\begin{lem}
\label{lem:three-levels}
Let $X$ be a random variable with law $P_\nu$. For every \(x\in J\) with \(\bar F_\nu(x)>0\),
\begin{alignat}{2}
\partial_\nu\log f_\nu(x) &=s_\nu(x)=K_\nu(x)-\E[K_\nu(X)], \label{eq:lr-kernel}\\
\partial_\nu\log\bar F_\nu(x) &=\E[s_\nu(X)\mid X\ge x]=\E[K_\nu(X)\mid X\ge x]-\E[K_\nu(X)], \label{eq:st-kernel}\\
\partial_\nu\log h_\nu(x) &=s_\nu(x)-\E[s_\nu(X)\mid X\ge x]=K_\nu(x)-\E[K_\nu(X)\mid X\ge x]. \label{eq:hr-kernel}
\end{alignat}
\end{lem}

\begin{equation}\label{eq:pair-density-int}
L_{\nu_2,\nu_1}(x) =\int_{\nu_1}^{\nu_2}s_\nu(x)\,\dx[\nu] =\int_{\nu_1}^{\nu_2} \bigl(K_\nu(x)-\E[K_\nu(X)]\bigr)\,\dx[\nu],
\end{equation}

\begin{thm}
\label{thm:four-criteria}
\leavevmode
\begin{enumerate}[label=\textup{(\roman*)}]
\item \(P_{\nu_1}\lr P_{\nu_2}\) for all \(\nu_1\le\nu_2\) in \(\mathcal I\) if and only if \(K_\nu\) is nondecreasing on \(J\) for every \(\nu\in\mathcal I\).

\item \(P_{\nu_2}\lc P_{\nu_1}\) for all \(\nu_1\le\nu_2\) in \(\mathcal I\) if and only if \(K_\nu\) is concave on \(J\) for every \(\nu\in\mathcal I\).

\item \(P_{\nu_1}\st P_{\nu_2}\) for all \(\nu_1\le\nu_2\) in \(\mathcal I\) if and only if \(\E[K_\nu(X)\mid X\ge x]\ge\E[K_\nu(X)]\) for every \(\nu\in\mathcal I\) and \(x\in J\) with \(\bar F_\nu(x)>0\).

\item \(P_{\nu_1}\hr P_{\nu_2}\) for all \(\nu_1\le\nu_2\) in \(\mathcal I\) if and only if \(K_\nu(x)\le\E[K_\nu(X)\mid X\ge x]\) for every \(\nu\in\mathcal I\) and \(x\in J\) with \(\bar F_\nu(x)>0\).
\end{enumerate}
\end{thm}

\begin{proof}
Sufficiency of \textup{(i)} and \textup{(ii)} follows from~\eqref{eq:pair-density-int}: the centring term is irrelevant, and an integral over \(\nu\) of nondecreasing or concave functions of \(x\) is again nondecreasing or concave. For \textup{(iii)} and \textup{(iv)}, integrate~\eqref{eq:st-kernel} or~\eqref{eq:hr-kernel} over $[\nu_1,\nu_2]$ and read off the sign.

For necessity,  identity~\eqref{def.log.companion} gives
\begin{equation}
\label{eq:kernel-score-shape}
K_\nu(x_1)-K_\nu(x_2)=s_\nu(x_1)-s_\nu(x_2),\qquad x_1,x_2\in J,
\end{equation}
and analogously for second differences and tail-conditional comparisons:
\begin{equation}
\label{eq:kernel-score-tail}
\E[K_\nu(X)\mid X\ge x]-\E[K_\nu(X)] =\E[s_\nu(X)\mid X\ge x]
\end{equation}
on \(\{x:\bar F_\nu(x)>0\}\), where the right-hand side equals \(\partial_\nu\log\bar F_\nu(x)\) by Lemma~\ref{lem:three-levels}. Both right-hand sides of (\ref{eq:kernel-score-shape}) and (\ref{eq:kernel-score-tail}) are continuous in \(\nu\) at every fixed \(x\) by Assumption~\ref{asm:diff-under-integral}. The kernel expressions on the left are therefore also continuous in \(\nu\), independently of the choice of $K_\nu$. The rest of the proof proceeds by contradiction, as follows:

\emph{Assertion~\textup{(i)}.} Suppose \(P_{\nu_1}\lr P_{\nu_2}\) for all \(\nu_1\le\nu_2\), and assume by contradiction that for some \(\nu_0 \in \mathcal I\), \(K_{\nu_0}\) is not nondecreasing on \(J\). By~\eqref{eq:kernel-score-shape}, this means $s_{\nu_0}$ is not nondecreasing, so there exist \(x_1<x_2\) in \(J\) such that \(s_{\nu_0}(x_1)>s_{\nu_0}(x_2)\). By continuity of $\nu\mapsto s_\nu(x_i)$, there is an \(\varepsilon>0\) and an interval $I_\varepsilon\subseteq\mathcal I$ containing $\nu_0$ such that $s_\nu(x_1)-s_\nu(x_2)\ge\varepsilon$ for all $\nu\in I_\varepsilon$. Choosing \(\nu_1<\nu_2\) in \(I_\varepsilon\) gives $L_{\nu_2,\nu_1}(x_1)-L_{\nu_2,\nu_1}(x_2)
=\int_{\nu_1}^{\nu_2}\bigl(s_\nu(x_1)-s_\nu(x_2)\bigr)\,\dx[\nu]
\ge \varepsilon(\nu_2-\nu_1)>0,
$
contradicting \(P_{\nu_1}\lr P_{\nu_2}\).

\emph{Assertion~\textup{(ii)}.} Suppose \(P_{\nu_2}\lc P_{\nu_1}\) for all \(\nu_1\le\nu_2\), and assume by contradiction that \(K_{\nu_0}\) is not concave on \(J\) for some \(\nu_0 \in \mathcal I\). Then $s_{\nu_0}$ is not concave by~\eqref{eq:kernel-score-shape} applied to second differences. In the continuous case, there exist \(x_1<x_2<x_3\) in \(J\) and \(\lambda\in(0,1)\) satisfying \(x_2=\lambda x_1+(1-\lambda)x_3\) such that
\(s_{\nu_0}(x_2) <\lambda s_{\nu_0}(x_1)+(1-\lambda)s_{\nu_0}(x_3).\)
In the discrete case, the same display holds for an admissible triplet with \(\lambda=1/2\).
By continuity in \(\nu\), there is an \(\varepsilon>0\) and an interval $I_\varepsilon$ containing $\nu_0$ such that \(\lambda s_\nu(x_1)+(1-\lambda)s_\nu(x_3)-s_\nu(x_2)\ge\varepsilon\) for all \(\nu\in I_\varepsilon\). Choosing \(\nu_1<\nu_2\) in \(I_\varepsilon\) yields
\begin{alignat*}{2}
\lambda L_{\nu_2,\nu_1}(x_1) +(1-\lambda)L_{\nu_2,\nu_1}(x_3) -L_{\nu_2,\nu_1}(x_2)
&=\int_{\nu_1}^{\nu_2} \bigl[\lambda s_\nu(x_1)+(1-\lambda)s_\nu(x_3)-s_\nu(x_2)\bigr]\,\dx[\nu]\\
&\ge \varepsilon(\nu_2-\nu_1)>0,
\end{alignat*}
contradicting concavity of \(L_{\nu_2,\nu_1}\) on \(J\).

\emph{Assertion~\textup{(iii)}.} Suppose \(P_{\nu_1}\st P_{\nu_2}\) for all \(\nu_1\le\nu_2\) and assume \(\E[K_{\nu_0}(X)\mid X\ge x_0]-\E[K_{\nu_0}(X)]<0\) at some \((\nu_0,x_0)\) with \(\bar F_{\nu_0}(x_0)>0\). By~\eqref{eq:kernel-score-tail} and Lemma~\ref{lem:three-levels}, this is the value of $\partial_\nu\log\bar F_\nu(x_0)$ at $\nu=\nu_0$, which is continuous in $\nu$. Hence there is an \(\varepsilon>0\) and an interval \(I_\varepsilon\) containing \(\nu_0\) such that $\partial_\nu\log\bar F_\nu(x_0)<-\varepsilon$ for every $\nu\in I_\varepsilon$. Choosing \(\nu_1<\nu_2\) in \(I_\varepsilon\) gives $T_{\nu_2,\nu_1}(x_0)<0$, contradicting \(P_{\nu_1}\st P_{\nu_2}\).

\emph{Assertion~\textup{(iv)}.} Suppose \(P_{\nu_1}\hr P_{\nu_2}\) for all \(\nu_1\le\nu_2\) and assume \(K_{\nu_0}(x_0)-\E[K_{\nu_0}(X)\mid X\ge x_0]>0\) at some \((\nu_0,x_0)\) with \(\bar F_{\nu_0}(x_0)>0\). By~\eqref{def.log.companion} and~\eqref{eq:kernel-score-tail}, this difference equals
\[
s_{\nu_0}(x_0)-\E[s_{\nu_0}(X)\mid X\ge x_0]=\partial_\nu\log h_\nu(x_0)\Big|_{\nu=\nu_0}
\]
by Lemma~\ref{lem:three-levels}, which is continuous in $\nu$. Hence there is an \(\varepsilon>0\) and an interval \(I_\varepsilon\) containing \(\nu_0\) such that $\partial_\nu\log h_\nu(x_0)>\varepsilon$ for every $\nu\in I_\varepsilon$. Choosing \(\nu_1<\nu_2\) in \(I_\varepsilon\) then contradicts $H_{\nu_2,\nu_1}(x_0)\le 0$. This completes the proof.
\end{proof}

\end{document}
